\subsection{Convolution of functions on a group}

We conserve the notations G, \(\beta\) , \(\chi\) , \(\chi'\) of No.~4.

Recall that if f is a complex function on G, the property of being locally \(\beta\) -integrable is independent of the choice of \(\beta\) . Let \(\mathcal{L}(\mathrm{G})\) be the set of functions having this property. If \(f \in \mathcal{L}(\mathrm{G})\) , \(g \in \mathcal{L}(\mathrm{G})\) , the relation

\[\text{«}f \cdot \beta\text{ and }g \cdot \beta\text{ are convolvable »}\]

is independent of the choice of \(\beta\) (§3, No.~1, Prop. 6). We shall then say that \(f\) and \(g\) are convolvable. By No.~1, \((f\cdot \beta)*(g\cdot \beta)\) is of the form \(h\cdot \beta\) with \(h\in \mathcal{L}(\mathbf{G})\) , \(h\) being determined up to locally \(\beta\) -negligible sets. We shall write \(h = f * ^ {\beta} g\) and we shall say that \(h\) is a convolution product of \(f\) and \(g\) relative to \(\beta\) . (One omits \(\beta\) when no confusion is possible.) If \(\beta\) is replaced by \(\psi \cdot \beta\) , \(\psi\) being a continuous representation of \(\mathbf{G}\) in \(\mathbf{R}_+^*\) , \(h\) does not change (§3, No.~1, Prop. 6); if \(\beta\) is replaced by \(a\beta\) ( \(a\in \mathbf{R}_+^*\) ), then \(h\) is replaced by ah. The convolution product of several functions on G is defined in an analogous manner.

If one of the convolutions of f and g is continuous, it is uniquely determined since the support of \(\beta\) is G. It is then called the convolution product of f and g relative to \(\beta\) .

It is clear that

\[
f * ^ {\beta} g = (f \cdot \beta) * ^ {\beta} g = f * ^ {\beta} (g \cdot \beta).\tag{14}
\]

PROPOSITION 9. --- Let \(f,g\) be in \(\mathcal{L}(\mathbf{G})\) . Assume that the function \(s\mapsto g(s^{-1}x)f(s)\chi(s^{-1})\) is essentially \(\beta\) -integrable except for a locally \(\beta\) -negligible set of values of \(x\) , and that the function

\[
x \mapsto \int | g (s ^ {- 1} x) f (s) | \chi (s ^ {- 1}) d \beta (s),
\]

defined locally \(\beta\) -almost everywhere, is locally \(\beta\) -integrable. Then \(f\) and \(g\) are convolvable.

This follows from Prop. 1 of No.~1.

PROPOSITION 10. --- Let \(f, g\) be in \(\mathcal{L}(G)\) . Assume that one of these two functions is continuous or is zero on the complement of a countable union of compact sets. If \(f\) and \(g\) are convolvable, then the function \(f * g\) is given locally \(\beta\) -almost everywhere by

\[
\begin{array}{c} (f * g) (x) = \int_ {\mathrm{G}} g (s ^ {- 1} x) f (s) \chi (s ^ {- 1}) d \beta (s) \\ = \int_ {\mathrm{G}} f (x s ^ {- 1}) g (s) \chi^ {\prime} (s ^ {- 1}) d \beta (s). \end{array}\tag{15}
\]

This follows from Prop. 2 of No.~1, and the remarks in No.~4.

In particular, if \(f * g\) is continuous and is given at every point by (15), then

\[
(f * g) (e) = \int g (s ^ {- 1}) f (s) \chi (s ^ {- 1}) d \beta (s) = \int f (s ^ {- 1}) g (s) \chi^ {\prime} (s ^ {- 1}) d \beta (s).\tag{16}
\]

Still more particularly, if \(\beta\) is a left and right Haar measure, and if \(f*g\) and \(g*f\) are continuous and are given at every point by (15) and the analogous formula for \(g*f\) , then

\[
(f * g) (e) = (g * f) (e) = \int f (s) g (s ^ {- 1}) d \beta (s).\tag{17}
\]

PROPOSITION 11. --- Let \(f,g\) be in \(\mathcal{L}(\mathrm{G})\) . Assume that one of the functions \(f,g\) is continuous, and that one of the functions \(f,g\) has compact support. Then \(f\) and \(g\) are convolvable. The formula (15) defines for all \(x\in\mathrm{G}\) a product \(f*g\) that is continuous. If \(f\in\mathcal{K}(\mathrm{G})\) and \(g\in\mathcal{K}(\mathrm{G})\) , then \(f*g\in\mathcal{K}(\mathrm{G})\) .

This follows from Props. 3 and 4 of No.~2.

PROPOSITION 12. --- Let \(p\) and \(q\) be two conjugate exponents \((1 \leqslant p \leqslant +\infty)\) . If \(f\chi^{-1/q} \in \mathrm{L}^1(\mathrm{G}, \beta)\) and \(g \in \mathrm{L}^p(\mathrm{G}, \beta)\) , then \(f\) and \(g\) are convolvable, \(f * g\) is equal locally \(\beta\) -almost everywhere to a function in \(\mathrm{L}^p(\mathrm{G}, \beta)\) , and

\[
\left\| f * g \right\| _ {p} \leqslant \left\| f \chi^ {- 1 / q} \right\| _ {1} \left\| g \right\| _ {p}.
\]

If \(f \in \mathrm{L}^p(\mathbf{G},\beta)\) and \(g\chi'^{-1/q} \in \mathrm{L}^1(\mathbf{G},\beta)\) , then \(f\) and \(g\) are convolvable, \(f * g\) is equal locally \(\beta\) -almost everywhere to a function in \(\mathrm{L}^p(\mathbf{G},\beta)\) , and

\[
\| f * g \| _ {p} \leqslant \| f \| _ {p} \| g \chi^ {\prime - 1 / q} \| _ {1}.
\]

This follows from Props. 5 and 6 of No.~2 and the remarks in No.~4.

PROPOSITION 13. --- If \(f\chi^{-1} \in \mathbf{L}^1(\mathbf{G}, \beta)\) and \(g \in \overline{\mathcal{H}(\mathbf{G})}\) , or if \(f \in \overline{\mathcal{H}(\mathbf{G})}\) and \(g\chi'^{-1} \in \mathbf{L}^1(\mathbf{G}, \beta)\) , then \(f\) and \(g\) are convolvable, and (15) defines for every \(x \in \mathbf{G}\) a product \(f * g\) that belongs to \(\overline{\mathcal{H}(\mathbf{G})}\) .

This follows from Prop. 5 of No.~2, and the remarks in No.~4.

PROPOSITION 14. --- If \(f\chi^{-1} \in \mathbf{L}^1(\mathbf{G}, \beta)\) and \(g \in \mathbf{L}^\infty(\mathbf{G}, \beta)\) , then the formula (15) defines for every \(x \in \mathbf{G}\) a product \(f * g\) that is bounded and is uniformly continuous for the right uniform structure of \(\mathbf{G}\) .

We already know that \(f * g\) belongs to \(\mathrm{L}^{\infty}(\mathrm{G},\beta)\) (No.~2, Prop. 5); moreover, \((f*g)(x) = \int f(xs^{-1})g(s)d\nu(s)\) , on setting \(\nu = \chi'^{-1}\cdot \beta\) ; \(\nu\) is a right Haar measure. Therefore

\[
\begin{array}{c} | (f * g) (x) - (f * g) (x ^ {\prime}) | \leqslant \| g \| _ {\infty} \int | f (x s ^ {- 1}) - f (x ^ {\prime} s ^ {- 1}) | d \nu (s) \\ = \| g \| _ {\infty} \int | (f (s ^ {- 1}) - f (x ^ {\prime} x ^ {- 1} s ^ {- 1}) d \nu (s) \end{array}
\]

and the latter integral is arbitrarily small provided \(x'x^{-1}\) is in a suitable neighborhood of \(e\) (§2, No.~5, Prop. 8).

PROPOSITION 15. --- Let \(p\) and \(q\) be two conjugate exponents \((1 < p < +\infty)\) . Assume that \(\beta\) is left-invariant. Let \(f \in \mathrm{L}^p(\mathrm{G}, \beta)\) , \(g \in \mathrm{L}^q(\mathrm{G}, \breve{\beta})\) . Then \(f\) and \(g\) are convolvable. The formula (15) defines, for every \(x \in \mathrm{G}\) , a product \(f * g\) that belongs to \(\overline{\mathcal{H}(\mathrm{G})}\) and is such that \(\| f * g \|_{\infty} \leqslant \| f \|_p \| \breve{g} \|_q\) .

For, we have \(\check{g} \in \mathbf{L}^q(\mathbf{G}, \beta)\) , therefore the function \(s \mapsto g(s^{-1}x)f(s)\) is \(\beta\) -integrable for every \(x \in \mathbf{G}\) . Moreover,

\[
\begin{array}{c} \int | g (s ^ {- 1} x) f (s) | d \beta (s) \leqslant \bigg (\int | f (s) | ^ {p} d \beta (s) \bigg) ^ {1 / p} \bigg (\int | g (s ^ {- 1} x) | ^ {q} d \beta (s) \bigg) ^ {1 / q} \\ = \| f \| _ {p} \bigg (\int | \check {g} (x ^ {- 1} s) | ^ {q} d \beta (s) \bigg) ^ {1 / q} = \| f \| _ {p} \| \check {g} \| _ {q}, \end{array}
\]

therefore \(f\) and \(g\) are convolvable (Prop. 9). One sees at the same time that (15) defines for every \(x\) a product \(f * g\) such that

\[
\left| (f * g) (x) \right| \leqslant \| f \| _ {p} \| \check {g} \| _ {q}.
\]

For \(f, g\) in \(\mathcal{H}(\mathrm{G})\) , we have \(f * g \in \mathcal{H}(\mathrm{G})\) (Prop. 11); therefore, for \(f \in \mathrm{L}^p(\mathrm{G}, \beta)\) and \(g \in \mathrm{L}^q(\mathrm{G}, \beta)\) , the product \(f * g\) furnished by (15) is the uniform limit of functions in \(\mathcal{H}(\mathrm{G})\) , hence belongs to \(\overline{\mathcal{H}(\mathrm{G})}\) .

COROLLARY. --- Let \(f \in \mathbf{L}^2(\mathbf{G}, \beta)\) , \(g \in \mathbf{L}^2(\mathbf{G}, \beta)\) . Then \(f\) and \(\breve{g}\) are convolvable. One of the convolutions \(f * \breve{g}\) belongs to \(\overline{\mathcal{H}(\mathbf{G})}\) and its value at \(e\) is \(\int_{\mathbf{G}} f(s) \overline{g(s)} d\beta(s)\) .

It suffices to take \(p = q = 2\) in Prop. 15 and to apply (16).

We no longer assume \(\beta\) to be left-invariant. Let \(\rho\) be a lower semi-continuous finite function \textgreater0 on G, such that \(\rho(st) \leqslant \rho(s)\rho(t)\) for all s, t in G. We denote by \(\mathrm{L}^{\rho}(\mathrm{G}, \beta)\) the set of equivalence classes of the complex functions on G that are integrable for \(\rho \cdot \beta\) . By the mapping \(f \mapsto f \cdot \beta\) , \(\mathrm{L}^{\rho}(\mathrm{G}, \beta)\) may be identified with the set of elements of \(\mathcal{M}^{\rho}(\mathrm{G})\) that have base \(\beta\) (a set that is independent of the choice of \(\beta\) ). If one sets

\[
\| f \| _ {\rho} = \int_ {G} | f (s) | \rho (s) d \beta (s)
\]

for \(f \in \mathrm{L}^{\rho}(\mathrm{G},\beta)\) , this identification is compatible with the norms, thus \(\mathrm{L}^{\rho}(\mathrm{G},\beta)\) appears as a complete normed subalgebra of \(\mathcal{M}^{\rho}(\mathrm{G})\) . It is even a two-sided ideal of \(\mathcal{M}^{\rho}(\mathrm{G})\) by Prop. 10 of §3, No.~2. (For \(\rho = 1\) , one recovers one of the assertions of No.~4.) In particular, \(\mathrm{L}^{1}(\mathrm{G},\beta)\) may be identified with a closed two-sided ideal of \(\mathcal{M}^{1}(\mathrm{G})\) .

PROPOSITION 16. --- Let U be a continuous representation of G in a Banach space E. Set \(\rho(s) = \|U(s)\|\) for all \(s \in G\) . For every \(f \in L^{\rho}(G, \beta)\) , set \(U(f) = U(f \cdot \beta)\) . Then \(f \mapsto U(f)\) is a linear representation of the algebra \(L^{\rho}(G, \beta)\) in E, such that \(\|U(f)\| \leqslant \|f\|_{\rho}\) .

This follows from §2, No.~6 and §3, No.~3, Prop. 11.

